% GENERATED FILE. Edit canonical lessons, then run npm run generate:books.
% canonical-source-hash: fnv1a64:b7ffb2c4a34c700c
% canonical-lessons: MW-C261-masi, MW-C261-bhatijo, MW-C261-bhanjo, MW-C261-poto, MW-C261-dohito

\chapter{Aunts, Nephews and Grandsons}
\label{ch:mw-aunts-nephews-and-grandsons}
\hlchaptermodality{261}

\begin{chapteropening}
\textbf{By the end of this chapter:} \emph{I can name my aunt, my nephews and my grandsons.}

Complete the last lesson of chapter 261: I can name my aunt, my nephews and my grandsons.
\end{chapteropening}

\section[\texorpdfstring{\mw{मासी} (māsī)}{मासी (māsī)}]{\mw{मासी} (māsī) — a mother's sister; an aunt}

\label{lesson:MW-C261-masi}



\begin{quote}
Before the new one: say the Marwadi for a sister's husband, then the Marwadi for a wife's brother.
\end{quote}

\subsection*{You'll want to know: \mw{मासी}}
\textbf{\mw{मासी}} — \emph{māsī} — \textquotedblleft{}a mother's sister; an aunt\textquotedblright{}.

\begin{quote}
\mw{म्हारी} \mw{मासी} \mw{गांव} \mw{मांय} \mw{रैवै} \mw{है।}
\end{quote}

\begin{quote}
\emph{mhārī māsī gā̃v mā̃y raivai hai.}
\end{quote}

\begin{quote}
My mother's sister lives in the village.
\end{quote}

A feminine noun.

\begin{cousinweb}
\textbf{\mw{मासी}} (\emph{māsī}) is inherited from Sanskrit \emph{mātuḥṣvasṛ}, \textquotedblleft{}a mother's sister\textquotedblright{}. CDIAL gives no Marwari form in this entry; it lists Hindi \emph{māsī} (CDIAL 10001).
\end{cousinweb}

\begin{grammarlens}[title={What you've built}]
One more word for this chapter.
\end{grammarlens}

\subsection*{Your turn}
\begin{itemize}
\raggedright
  \item \emph{Say it:} \emph{māsī}
  \item \emph{Say it:} \emph{māsī}, once more
  \item \emph{Say it:} \emph{mhārī māsī gā̃v mā̃y raivai hai.} — My mother's sister lives in the village.
  \item \emph{Recall:} say the Marwadi for a sister's husband, then the Marwadi for a wife's brother, then say \emph{māsī} again
\end{itemize}

\begin{tcolorbox}[breakable,colback=teal!4,colframe=teal!35!black,title={Before you move on}]
What does \mw{मासी} mean? (\textquotedblleft{}A mother's sister; an aunt\textquotedblright{}.) Say it once more.
\end{tcolorbox}
\section[\texorpdfstring{\mw{भतीजो} (bhatījo)}{भतीजो (bhatījo)}]{\mw{भतीजो} (bhatījo) — a brother's son; a nephew}

\label{lesson:MW-C261-bhatijo}



\begin{quote}
Before the new one: say the Marwadi for a wife's brother, then the Marwadi for a mother's sister.
\end{quote}

\subsection*{You'll want to know: \mw{भतीजो}}
\textbf{\mw{भतीजो}} — \emph{bhatījo} — \textquotedblleft{}a brother's son; a nephew\textquotedblright{}.

\begin{quote}
\mw{म्हारो} \mw{भतीजो} \mw{काल} \mw{आयो।}
\end{quote}

\begin{quote}
\emph{mhāro bhatījo kāl āyo.}
\end{quote}

\begin{quote}
My brother's son came yesterday.
\end{quote}

A masculine noun.

\begin{cousinweb}
\textbf{\mw{भतीजो}} (\emph{bhatījo}) is inherited from Sanskrit \emph{bhrātrīya}, \textquotedblleft{}a brother's son\textquotedblright{}. CDIAL gives no Marwari form in this entry; it lists Hindi \emph{bhatījā} (CDIAL 9672).
\end{cousinweb}

\begin{grammarlens}[title={What you've built}]
One more word for this chapter.
\end{grammarlens}

\subsection*{Your turn}
\begin{itemize}
\raggedright
  \item \emph{Say it:} \emph{bhatījo}
  \item \emph{Say it:} \emph{bhatījo}, once more
  \item \emph{Say it:} \emph{mhāro bhatījo kāl āyo.} — My brother's son came yesterday.
  \item \emph{Recall:} say the Marwadi for a wife's brother, then the Marwadi for a mother's sister, then say \emph{bhatījo} again
\end{itemize}

\begin{tcolorbox}[breakable,colback=teal!4,colframe=teal!35!black,title={Before you move on}]
What does \mw{भतीजो} mean? (\textquotedblleft{}A brother's son; a nephew\textquotedblright{}.) Say it once more.
\end{tcolorbox}
\section[\texorpdfstring{\mw{भाणजो} (bhāṇjo)}{भाणजो (bhāṇjo)}]{\mw{भाणजो} (bhāṇjo) — a sister's son; a nephew}

\label{lesson:MW-C261-bhanjo}



\begin{quote}
Before the new one: say the Marwadi for a mother's sister, then the Marwadi for a brother's son.
\end{quote}

\subsection*{You'll want to know: \mw{भाणजो}}
\textbf{\mw{भाणजो}} — \emph{bhāṇjo} — \textquotedblleft{}a sister's son; a nephew\textquotedblright{}.

\begin{quote}
\mw{म्हारो} \mw{भाणजो} \mw{सहर} \mw{सूं} \mw{आयो।}
\end{quote}

\begin{quote}
\emph{mhāro bhāṇjo sahar sū̃ āyo.}
\end{quote}

\begin{quote}
My sister's son came from the city.
\end{quote}

A masculine noun.

\begin{cousinweb}
\textbf{\mw{भाणजो}} (\emph{bhāṇjo}) is inherited from Sanskrit \emph{bhāgineya}, \textquotedblleft{}a sister's son\textquotedblright{}. CDIAL gives no Marwari form in this entry; it lists Gujarati \emph{bhāṇej} and \emph{bhāṇiyo} (CDIAL 9433).
\end{cousinweb}

\begin{grammarlens}[title={What you've built}]
One more word for this chapter.
\end{grammarlens}

\subsection*{Your turn}
\begin{itemize}
\raggedright
  \item \emph{Say it:} \emph{bhāṇjo}
  \item \emph{Say it:} \emph{bhāṇjo}, once more
  \item \emph{Say it:} \emph{mhāro bhāṇjo sahar sū̃ āyo.} — My sister's son came from the city.
  \item \emph{Recall:} say the Marwadi for a mother's sister, then the Marwadi for a brother's son, then say \emph{bhāṇjo} again
\end{itemize}

\begin{tcolorbox}[breakable,colback=teal!4,colframe=teal!35!black,title={Before you move on}]
What does \mw{भाणजो} mean? (\textquotedblleft{}A sister's son; a nephew\textquotedblright{}.) Say it once more.
\end{tcolorbox}
\section[\texorpdfstring{\mw{पोतो} (poto)}{पोतो (poto)}]{\mw{पोतो} (poto) — a son's son; a grandson}

\label{lesson:MW-C261-poto}



\begin{quote}
Before the new one: say the Marwadi for a brother's son, then the Marwadi for a sister's son.
\end{quote}

\subsection*{You'll want to know: \mw{पोतो}}
\textbf{\mw{पोतो}} — \emph{poto} — \textquotedblleft{}a son's son; a grandson\textquotedblright{}.

\begin{quote}
\mw{म्हारो} \mw{पोतो} \mw{घणो} \mw{छोटो} \mw{है।}
\end{quote}

\begin{quote}
\emph{mhāro poto ghaṇo chhoṭo hai.}
\end{quote}

\begin{quote}
My grandson is very small.
\end{quote}

A masculine noun.

\begin{cousinweb}
\textbf{\mw{पोतो}} (\emph{poto}) is inherited from Sanskrit \emph{pautra}, \textquotedblleft{}a son's son\textquotedblright{}. CDIAL gives no Marwari form in this entry; it lists Hindi \emph{potā} (CDIAL 8416).
\end{cousinweb}

\begin{grammarlens}[title={What you've built}]
One more word for this chapter.
\end{grammarlens}

\subsection*{Your turn}
\begin{itemize}
\raggedright
  \item \emph{Say it:} \emph{poto}
  \item \emph{Say it:} \emph{poto}, once more
  \item \emph{Say it:} \emph{mhāro poto ghaṇo chhoṭo hai.} — My grandson is very small.
  \item \emph{Recall:} say the Marwadi for a brother's son, then the Marwadi for a sister's son, then say \emph{poto} again
\end{itemize}

\begin{tcolorbox}[breakable,colback=teal!4,colframe=teal!35!black,title={Before you move on}]
What does \mw{पोतो} mean? (\textquotedblleft{}A son's son; a grandson\textquotedblright{}.) Say it once more.
\end{tcolorbox}
\section[\texorpdfstring{\mw{दोहितो} (dohito)}{दोहितो (dohito)}]{\mw{दोहितो} (dohito) — a daughter's son; a grandson}

\label{lesson:MW-C261-dohito}



\begin{quote}
Before the new one: say the Marwadi for a sister's son, then the Marwadi for a son's son.
\end{quote}

\subsection*{You'll want to know: \mw{दोहितो}}
\textbf{\mw{दोहितो}} — \emph{dohito} — \textquotedblleft{}a daughter's son; a grandson\textquotedblright{}.

\begin{quote}
\mw{म्हारो} \mw{दोहितो} \mw{गीत} \mw{गावै} \mw{है।}
\end{quote}

\begin{quote}
\emph{mhāro dohito gīt gāvai hai.}
\end{quote}

\begin{quote}
My daughter's son sings songs.
\end{quote}

A masculine noun.

\begin{cousinweb}
\textbf{\mw{दोहितो}} (\emph{dohito}) is inherited from Sanskrit \emph{dauhitra}, \textquotedblleft{}a daughter's son\textquotedblright{}. CDIAL gives no Marwari form in this entry; it lists Hindi \emph{dohtā} and Gujarati \emph{dohitra} (CDIAL 6605).
\end{cousinweb}

\begin{grammarlens}[title={What you've built}]
One more word for this chapter.
\end{grammarlens}

\subsection*{Your turn}
\begin{itemize}
\raggedright
  \item \emph{Say it:} \emph{dohito}
  \item \emph{Say it:} \emph{dohito}, once more
  \item \emph{Say it:} \emph{mhāro dohito gīt gāvai hai.} — My daughter's son sings songs.
  \item \emph{Recall:} say the Marwadi for a sister's son, then the Marwadi for a son's son, then say \emph{dohito} again
\end{itemize}

\begin{tcolorbox}[breakable,colback=teal!4,colframe=teal!35!black,title={Before you move on}]
What does \mw{दोहितो} mean? (\textquotedblleft{}A daughter's son; a grandson\textquotedblright{}.) Say it once more.
\end{tcolorbox}
